% GENERATED FILE. Edit canonical lessons, then run npm run generate:books.
% canonical-source-hash: fnv1a64:f20b14c3533d5b59
% canonical-lessons: MR-C358-samvidhan, MR-C358-dhoran, MR-C358-lagu-karne, MR-C358-radda-karne, MR-C358-bahal-karne

\chapter{Making Laws}
\label{ch:mr-making-laws}
\hlchaptermodality{358}

\begin{chapteropening}
\textbf{By the end of this chapter:} \emph{I can talk about laws and policies being made, changed or cancelled.}

Complete the last lesson of chapter 358: I can talk about laws and policies being made, changed or cancelled.
\end{chapteropening}

\section[\texorpdfstring{\mr{संविधान} (saṁvidhān)}{संविधान (saṁvidhān)}]{\mr{संविधान} (saṁvidhān) — a constitution}

\label{lesson:MR-C358-samvidhan}



\begin{quote}
Before the new one: say the Marathi for a leader, then the Marathi for to appoint.
\end{quote}

\subsection*{You'll want to know: \mr{संविधान}}
\textbf{\mr{संविधान}} — \emph{saṁvidhān} — \textquotedblleft{}a constitution\textquotedblright{}.

\begin{quote}
\mr{संविधानाने} \mr{प्रत्येक} \mr{नागरिकाला} \mr{समान} \mr{हक्क} \mr{दिले} \mr{आहेत}.
\end{quote}

\begin{quote}
\emph{saṁvidhānāne pratyek nāgarikālā samān hakka dile āhet.}
\end{quote}

\begin{quote}
The constitution has given every citizen equal rights.
\end{quote}

Neuter.

Neuter; also \mr{राज्यघटना}.

\begin{cousinweb}
\textbf{\mr{संविधान}} (\emph{saṁvidhān}) is a tatsama, taken from Sanskrit \emph{saṁvidhāna}, \textquotedblleft{}arrangement, disposition\textquotedblright{} (Monier-Williams).
\end{cousinweb}

\begin{grammarlens}[title={What you've built}]
One more word for giving reasons and opinions.
\end{grammarlens}

\subsection*{Your turn}
\begin{itemize}
\raggedright
  \item \emph{Say it:} \emph{saṁvidhān}
  \item \emph{Say it:} \emph{saṁvidhān}, once more
  \item \emph{Say it:} \emph{saṁvidhānāne pratyek nāgarikālā samān hakka dile āhet.} — The constitution has given every citizen equal rights.
  \item \emph{Recall:} say the Marathi for a leader, then the Marathi for to appoint, then say \emph{saṁvidhān} again
\end{itemize}

\begin{tcolorbox}[breakable,colback=teal!4,colframe=teal!35!black,title={Before you move on}]
What does \mr{संविधान} mean? (\textquotedblleft{}A constitution\textquotedblright{}.) Say it once more.
\end{tcolorbox}
\section[\texorpdfstring{\mr{धोरण} (dhoraṇ)}{धोरण (dhoraṇ)}]{\mr{धोरण} (dhoraṇ) — a policy}

\label{lesson:MR-C358-dhoran}



\begin{quote}
Before the new one: say the Marathi for to appoint, then the Marathi for a constitution.
\end{quote}

\subsection*{You'll want to know: \mr{धोरण}}
\textbf{\mr{धोरण}} — \emph{dhoraṇ} — \textquotedblleft{}a policy\textquotedblright{}.

\begin{quote}
\mr{सरकारचं} \mr{नवं} \mr{शिक्षण} \mr{धोरण} \mr{जाहीर} \mr{झालं}.
\end{quote}

\begin{quote}
\emph{sarkārcaṁ navaṁ śikṣaṇ dhoraṇ jāhīr jhālaṁ.}
\end{quote}

\begin{quote}
The government's new education policy was announced.
\end{quote}

Neuter.

Neuter: \mr{नवं} \mr{धोरण}.

\begin{cousinweb}
\textbf{\mr{धोरण}} (\emph{dhoraṇ}) is a Marathi noun whose origin is not settled.
\end{cousinweb}

\begin{grammarlens}[title={What you've built}]
One more word for giving reasons and opinions.
\end{grammarlens}

\subsection*{Your turn}
\begin{itemize}
\raggedright
  \item \emph{Say it:} \emph{dhoraṇ}
  \item \emph{Say it:} \emph{dhoraṇ}, once more
  \item \emph{Say it:} \emph{sarkārcaṁ navaṁ śikṣaṇ dhoraṇ jāhīr jhālaṁ.} — The government's new education policy was announced.
  \item \emph{Recall:} say the Marathi for to appoint, then the Marathi for a constitution, then say \emph{dhoraṇ} again
\end{itemize}

\begin{tcolorbox}[breakable,colback=teal!4,colframe=teal!35!black,title={Before you move on}]
What does \mr{धोरण} mean? (\textquotedblleft{}A policy\textquotedblright{}.) Say it once more.
\end{tcolorbox}
\section[\texorpdfstring{\mr{लागू} \mr{करणे} (lāgū karṇe)}{लागू करणे (lāgū karṇe)}]{\mr{लागू} \mr{करणे} (lāgū karṇe) — to enforce; to bring into force}

\label{lesson:MR-C358-lagu-karne}



\begin{quote}
Before the new one: say the Marathi for a constitution, then the Marathi for a policy.
\end{quote}

\subsection*{You'll want to know: \mr{लागू} \mr{करणे}}
\textbf{\mr{लागू} \mr{करणे}} — \emph{lāgū karṇe} — \textquotedblleft{}to enforce; to bring into force\textquotedblright{}.

\begin{quote}
\mr{सरकारने} \mr{नवा} \mr{नियम} \mr{लागू} \mr{केला}.
\end{quote}

\begin{quote}
\emph{sarkārne navā niyam lāgū kelā.}
\end{quote}

\begin{quote}
The government brought a new rule into force.
\end{quote}

Past, with -\mr{ने} (he …): \mr{त्याने} \mr{नियम} \mr{लागू} \mr{केला}, \emph{tyāne niyam lāgū kelā}.

\mr{लागू} \mr{होणे}, to apply: \mr{हा} \mr{नियम} \mr{सगळ्यांना} \mr{लागू} \mr{होतो}.

\begin{cousinweb}
\textbf{\mr{लागू} \mr{करणे}} (\emph{lāgū karṇe}) joins \mr{लागू}, \textquotedblleft{}in force, applying\textquotedblright{}, built on \mr{लागणे}, \textquotedblleft{}to attach\textquotedblright{}, which is inherited from Sanskrit \emph{lagyati}, \textquotedblleft{}adheres to\textquotedblright{} (CDIAL 10895), to \mr{करणे}.
\end{cousinweb}

\begin{grammarlens}[title={What you've built}]
One more word for giving reasons and opinions.
\end{grammarlens}

\subsection*{Your turn}
\begin{itemize}
\raggedright
  \item \emph{Say it:} \emph{lāgū karṇe}
  \item \emph{Say it:} \emph{lāgū karṇe}, once more
  \item \emph{Say it:} \emph{sarkārne navā niyam lāgū kelā.} — The government brought a new rule into force.
  \item \emph{Recall:} say the Marathi for a constitution, then the Marathi for a policy, then say \emph{lāgū karṇe} again
\end{itemize}

\begin{tcolorbox}[breakable,colback=teal!4,colframe=teal!35!black,title={Before you move on}]
What does \mr{लागू} \mr{करणे} mean? (\textquotedblleft{}To enforce; to bring into force\textquotedblright{}.) Say it once more.
\end{tcolorbox}
\section[\texorpdfstring{\mr{रद्द} \mr{करणे} (radda karṇe)}{रद्द करणे (radda karṇe)}]{\mr{रद्द} \mr{करणे} (radda karṇe) — to cancel; to repeal}

\label{lesson:MR-C358-radda-karne}



\begin{quote}
Before the new one: say the Marathi for a policy, then the Marathi for to enforce.
\end{quote}

\subsection*{You'll want to know: \mr{रद्द} \mr{करणे}}
\textbf{\mr{रद्द} \mr{करणे}} — \emph{radda karṇe} — \textquotedblleft{}to cancel; to repeal\textquotedblright{}.

\begin{quote}
\mr{पावसामुळे} \mr{शाळेने} \mr{सहल} \mr{रद्द} \mr{केली}.
\end{quote}

\begin{quote}
\emph{pāvsāmuḷe śāḷene sahal radda kelī.}
\end{quote}

\begin{quote}
Because of the rain, the school cancelled the trip.
\end{quote}

Past, with -\mr{ने} (he …): \mr{त्याने} \mr{रद्द} \mr{केलं}, \emph{tyāne radda kelaṁ}.

\mr{रद्द} \mr{होणे}, to be cancelled.

\begin{cousinweb}
\textbf{\mr{रद्द} \mr{करणे}} (\emph{radda karṇe}) joins \mr{रद्द}, \textquotedblleft{}void\textquotedblright{}, a loan from Persian usually traced to Arabic \emph{radd}, \textquotedblleft{}returning, rejection\textquotedblright{}, to \mr{करणे}, \textquotedblleft{}to make\textquotedblright{}.
\end{cousinweb}

\begin{grammarlens}[title={What you've built}]
One more word for giving reasons and opinions.
\end{grammarlens}

\subsection*{Your turn}
\begin{itemize}
\raggedright
  \item \emph{Say it:} \emph{radda karṇe}
  \item \emph{Say it:} \emph{radda karṇe}, once more
  \item \emph{Say it:} \emph{pāvsāmuḷe śāḷene sahal radda kelī.} — Because of the rain, the school cancelled the trip.
  \item \emph{Recall:} say the Marathi for a policy, then the Marathi for to enforce, then say \emph{radda karṇe} again
\end{itemize}

\begin{tcolorbox}[breakable,colback=teal!4,colframe=teal!35!black,title={Before you move on}]
What does \mr{रद्द} \mr{करणे} mean? (\textquotedblleft{}To cancel; to repeal\textquotedblright{}.) Say it once more.
\end{tcolorbox}
\section[\texorpdfstring{\mr{बहाल} \mr{करणे} (bahāl karṇe)}{बहाल करणे (bahāl karṇe)}]{\mr{बहाल} \mr{करणे} (bahāl karṇe) — to grant; to confer}

\label{lesson:MR-C358-bahal-karne}



\begin{quote}
Before the new one: say the Marathi for to enforce, then the Marathi for to cancel.
\end{quote}

\subsection*{You'll want to know: \mr{बहाल} \mr{करणे}}
\textbf{\mr{बहाल} \mr{करणे}} — \emph{bahāl karṇe} — \textquotedblleft{}to grant; to confer\textquotedblright{}.

\begin{quote}
\mr{संविधानाने} \mr{स्त्रियांना} \mr{मताचा} \mr{हक्क} \mr{बहाल} \mr{केला}.
\end{quote}

\begin{quote}
\emph{saṁvidhānāne striyānnā matācā hakka bahāl kelā.}
\end{quote}

\begin{quote}
The constitution granted women the right to vote.
\end{quote}

Past, with -\mr{ने} (he …): \mr{त्याने} \mr{बहाल} \mr{केलं}, \emph{tyāne bahāl kelaṁ}.

The receiver takes -\mr{ला}: \mr{स्त्रियांना} \mr{हक्क} \mr{बहाल} \mr{करणे}.

\begin{cousinweb}
\textbf{\mr{बहाल} \mr{करणे}} (\emph{bahāl karṇe}) joins \mr{बहाल}, \textquotedblleft{}confirmed, granted\textquotedblright{}, a loan usually traced to Persian \emph{ba-ḥāl}, \textquotedblleft{}in (its proper) state\textquotedblright{}, built on Arabic \emph{ḥāl}, \textquotedblleft{}state\textquotedblright{}, to \mr{करणे}, \textquotedblleft{}to make\textquotedblright{}.
\end{cousinweb}

\begin{grammarlens}[title={What you've built}]
One more word for giving reasons and opinions.
\end{grammarlens}

\subsection*{Your turn}
\begin{itemize}
\raggedright
  \item \emph{Say it:} \emph{bahāl karṇe}
  \item \emph{Say it:} \emph{bahāl karṇe}, once more
  \item \emph{Say it:} \emph{saṁvidhānāne striyānnā matācā hakka bahāl kelā.} — The constitution granted women the right to vote.
  \item \emph{Recall:} say the Marathi for to enforce, then the Marathi for to cancel, then say \emph{bahāl karṇe} again
\end{itemize}

\begin{tcolorbox}[breakable,colback=teal!4,colframe=teal!35!black,title={Before you move on}]
What does \mr{बहाल} \mr{करणे} mean? (\textquotedblleft{}To grant; to confer\textquotedblright{}.) Say it once more.
\end{tcolorbox}
